\section{Discussion}
\label{sec:discussion}

\subsection{What the Pilot Can and Cannot Conclude}

\todopilot{Open with one paragraph summarising the headline claim:
e.g., ``At 10M environment steps and three seeds per treatment, the
\emph{X} observation treatment achieved an IQM of \emph{Y}, with a
95\% bootstrap CI \emph{Z}; under the pre-registered probability-of-
improvement threshold, this difference is/is not statistically
significant.''  Be precise about what is and is not concluded.}

The 10M-step budget and three-seed pilot is deliberately preliminary.
It is sufficient to detect large effect sizes---a treatment that
clearly cannot leave Pallet Town versus one that consistently reaches
Viridian Forest---but it cannot reliably distinguish small-to-medium
effects.  In particular, the bootstrap CIs for the three treatments
\todopilot{do/do not} overlap, and we therefore \todopilot{can/cannot}
make a strong claim about the ranking of treatments at the
journal-version compute budget.

\subsection{Hypothesis 1 Revisited}
\label{sec:h1}

The pre-registered first hypothesis (H1) was that \emph{at matched
compute and matched architecture, symbolic observations yield
significantly better sample efficiency than pixel observations}.

\todopilot{Address H1 directly.  Did the data support, refute, or
fail-to-distinguish the hypothesis?  Mention the specific comparison
(IQM-symbolic vs IQM-pixel + their CIs) and the corresponding
probability-of-improvement value.  Be honest if the pilot is
inconclusive.}

\subsection{Why Pixel Might Underperform}

A pixel agent must implicitly re-derive all the structured state that
the symbolic agent reads directly from
RAM~\citep{anand2019unsupervised}: the player's coordinates, the
event-flag bit-vector, the contents of its party.  At a fixed compute
budget that re-derivation cost is paid in policy-network capacity and
in sample efficiency.  In a long-horizon, sparse-reward environment
where reward is gated on event flags, the symbolic agent has direct
access to the very signal that determines reward.  We expect this gap
to be especially pronounced early in training, when the pixel agent's
internal representations are not yet useful.

\subsection{Why Hybrid Might Not Strictly Dominate}

A naïve expectation is that the hybrid treatment, having access to
\emph{both} streams, should always win.  In practice, with a fixed
training budget, the hybrid agent must learn to weight two encoders
and may suffer from gradient interference if the two streams provide
redundant or conflicting signals.  Whether hybrid actually wins, ties,
or loses against the better single-stream treatment is an empirical
question; \todopilot{report the observed ranking}.

\subsection{Threats to Validity}

\paragraph{Encoder capacity is not parameter-matched.}
The hybrid encoder has roughly twice the parameter count of pixel or
symbolic alone.  This is a known confound; we discuss the trade-off in
\Cref{sec:obs-treatments}.  A capacity-matched ablation, in which the
single-stream encoders are widened to match the hybrid encoder's
total parameter count, is part of the journal-version plan.

\paragraph{Single fixed save state.}
All training in this pilot resets to \texttt{post\_intro.state}
(Pallet Town).  Effects observed here may not generalise to a
curriculum that samples from multiple save states, e.g., one rooted
inside a more partially-observable region of the game such as Mt.
Moon.  We will address this in the journal version with a
Go-Explore-style weighted save-state curriculum.

\paragraph{Three seeds.}
Three seeds is enough to detect large effects but provides limited
power for medium effects, particularly at the IQM tails.  All
quantitative claims are reported with bootstrap CIs and a
probability-of-improvement threshold; we do not over-interpret point
estimates.

\paragraph{No hyperparameter tuning per treatment.}
We use the same RecurrentPPO hyperparameters across all three
treatments, which is fair but pessimistic for a treatment whose
optimum lies far from those defaults.  The journal version will
include per-treatment hyperparameter tuning via Optuna on a held-out
save state.

\subsection{Future Work}
\label{sec:future-work}

The journal-length follow-up will extend this pilot in three
directions.  First, the seed count will increase to five (minimum)
and the training budget to $10^8$ steps per seed, restoring the
pre-registered design.  Second, four pre-committed
ablations---removing the LSTM, removing RND-style intrinsic
exploration, ablating frame-stacking, and ablating the save-state
curriculum---will isolate the contributions of recurrence, novelty
bonuses, temporal observation, and curriculum to the per-treatment
gap.  Third, a representational similarity analysis using
CKA~\citep{kornblith2019cka} will ask whether the three treatments
converge to similar internal features once trained, or whether the
performance gap reflects a deeper structural difference in what each
observation stream lets the agent represent.

\subsection{Open Release}

The full training, evaluation, and analysis pipeline is open-sourced.
The environment is pip-installable as \textsc{PokeGym}.  The W\&B
project is public and the run-level logs and model checkpoints are
linked from the supplementary material.  We hope the benchmark and
the analysis scaffolding lower the barrier for further work in
long-horizon RL on real, hand-designed game environments.
